% Generated draft: requires mathematical and source review.
\documentclass[11pt]{article}
\usepackage[margin=25mm]{geometry}
\usepackage{fontspec}
\setmainfont{Latin Modern Roman}
\usepackage{amsmath,amssymb,mathtools,unicode-math}
\setmathfont{Latin Modern Math}
\usepackage{enumitem,needspace,xurl,hyperref}
\setlength{\parindent}{0pt}
\setlength{\parskip}{5pt}
\setlength{\emergencystretch}{3em}
\newcommand{\N}{\mathbb{N}}
\newcommand{\Z}{\mathbb{Z}}
\newcommand{\Q}{\mathbb{Q}}
\newcommand{\R}{\mathbb{R}}
\newcommand{\C}{\mathbb{C}}
\newcommand{\F}{\mathbb{F}}
\newcommand{\A}{\mathbb{A}}
\newcommand{\PP}{\mathbb{P}}
\newcommand{\E}{\mathbb{E}}
\newcommand{\eps}{\varepsilon}
\newcommand{\dd}{\,\mathrm{d}}
\begin{document}
\section*{UnsolvedMath ID: 30005241}
\textbf{Source problem code:} catalogue entry
\subsection{Definitions and mathematical statement}
The supplied record concerns multiple scattering in wave scattering problems. For bookkeeping only, write a multiple-scattering expansion formally as
\[
\sum_{n=0}^{\infty} T_n,
\]
where \(T_n\) denotes the aggregate contribution at scattering order \(n\). This notation does not assert convergence or specify the space containing the contributions. For a truncation order \(N\), the expression \(\sum_{n>N}T_n\) denotes a formal tail; if the series diverges, it does not automatically define a remainder.

\textbf{Source-stated research objective.} Develop a collective method for computing the exponentially many single-scattering returns arising in the multiple-scattering construction, and a mathematically justified treatment of the remaining infinite tail, including cases where the multiple-scattering series diverges. The retained passage suggests stabilization of multiple-scattering ray paths as a possible way to address the preceding computational issue, without asserting that it succeeds; it identifies the infinite tail as an additional difficulty.

\textbf{Unresolved formulation.} The supplied excerpt begins mid-discussion and does not define the scattering problem, single-scattering returns, ray-path stabilization, or the expansion terms. It also does not identify the parameter governing exponential growth, the desired computed quantity, the relevant function spaces or norms, admissible configurations, or accuracy and complexity requirements. In the divergent case, no summation prescription or independently defined quantity to be approximated is specified. Consequently, the record supports this research objective but does not determine a precise quantified existence or error-estimate statement. These missing choices require source clarification and mathematical review.
\subsection{Short English statement}
The problem asks for a collective computation of exponentially many scattering returns and a justified treatment of the infinite multiple-scattering tail, even when the series diverges. The supplied passage suggests possible stabilization of ray paths but leaves the scattering model, accuracy requirements, and meaning of a divergent tail unspecified.
\subsection{Sources}
\begin{itemize}
\item Supplied local UnsolvedMath registry record, permanent ID \(30005241\), title \emph{Infinite Tails of Multiple-Scattering Series}, exact source code \texttt{catalogue entry}. The record assigns proposed year 2022 and supplies no individual proposer.
\item Supplied source citation: \emph{At the Interface between Semiclassical Analysis and Numerical Analysis of Wave Scattering Problems} (2023), Oberwolfach Reports, pp.~2511--2587, DOI \texttt{10.4171/owr/2022/43}. Supplied source and literature link: \url{https://doi.org/10.4171/owr/2022/43}, labelled in the literature sources as ``Oberwolfach Report 43/2022.'' The background describes these reports as records of mathematical workshops and their open-problem discussions.
\item Literature assessment dated exactly \texttt{2026-08-22}: the supplied curation labels the problem open and states that collective computation of exponentially many returns and treatment of the infinite tail remain unresolved, with possible divergence of the series. This is supplied dated curation, not an independently verified literature assessment; the background expressly calls for independent verification of time-sensitive status and cited claims.
\item The registry labels the statement \texttt{corrected\_verified} and says that it was reconstructed and checked against the official report. These are supplied verification claims, not verification performed for this draft. The retained original passage begins mid-discussion; the collective-computation objective is stated explicitly in the reconstructed statement.
\end{itemize}
\end{document}
